% Chapter 7 -- The Move Command (AN-08, AN-09).
% Sources (Law 1): os/i123/app/face_a.c menu_select / range_accept (the
% FROM and TO prompts; the single-cell TO rule), os/i123/core/range.c
% i123_transfer mode MOVE (the FROM range erased; formulas anywhere that
% refer into the FROM range re-encoded to follow it; range names whose
% corners lie in it moved; protected cells refuse), the copy-move
% reference rules graded against the real Release 2.2 MOVE01 goldens;
% spec/i123/era_shortfalls.md I123-P3-07.
% GATES: test-i123-copy-move-model test-i123-range-menu test-i123-menu-tree test-pubs-i123 test-pubs-i123-examples test-pubs-screens
\Gates{test-i123-copy-move-model, test-i123-range-menu, test-i123-menu-tree,
  test-pubs-i123, test-pubs-i123-examples, test-pubs-screens}
\chapter[The Move Command]{\textit{The Move Command}}
\input{i123-tree-Move}

The Move command moves the contents of a cell or a range to another part of
the worksheet.  Formulas that refer to the moved cells are adjusted so that
they still refer to the same data in its new place.

\begin{LotTable}{p{30mm}p{98mm}}
\LotHead{\textbf{Command} & \textbf{Task}}
/Move & Moves a cell or a range to another location in the worksheet\LotEnd
\end{LotTable}

\section{Common Uses of the Move Command}
\begin{itemize}
\item Rearranging the parts of a worksheet (\slashcmd{/Move})
\item Making room for new rows or columns of data in a table
  (\slashcmd{/Move})
\end{itemize}

\section{Related Reading}
\begin{itemize}
\item ``Formulas'' and ``Ranges'' in Chapter~1
\item Chapter~3, ``The Copy Command''
\item Lesson~5 in the \doctitle{Tutorial}
\end{itemize}

\section{Terms You Need to Know}
\begin{itemize}
\item The \term{FROM range} is the cell or range you move.
\item The \term{TO range} is the place it goes, named by its upper left
  cell.
\end{itemize}

\MenuEntry{/Move}
/Move moves a cell or a range, with the formats and protection settings of
its cells, to another part of the worksheet.

Initech 123 moves every cell of the FROM range to the cell in the same
position relative to the upper left cell of the TO range, and leaves the
FROM range empty.  Whatever the TO range held is replaced.  Formulas are
adjusted in two ways:
\begin{itemize}
\item a formula anywhere in the worksheet that refers to a cell of the FROM
  range now refers to that cell in its new place;
\item a moved formula's references to cells outside the FROM range still
  refer to the same cells.
\end{itemize}
A range name whose range lies in the FROM range moves with it.

\BHead{Procedure}
\begin{Steps123}
\item Select /Move.  The control panel shows
  \texttt{Enter range to move FROM:} followed by the address of the cell
  pointer.
\item Specify the FROM range and press \kname{ENTER}.
\item Specify the TO range, its upper left cell, and press \kname{ENTER}.
\end{Steps123}

\FlairShot[fig:movedone]{movedone}{B1..B3 moved to D1: the total in D3
still adds the two numbers above it}

\begin{Caution}
/Move replaces whatever the TO range contains.  A formula elsewhere that
referred to a cell of the TO range does not refer to the moved data.
\end{Caution}

\begin{Note}
Give the TO range as a single cell; a larger TO range displays
\texttt{Specify the upper left destination cell} and moves nothing.  If the
worksheet is protected and a cell of the FROM or the TO range is
protected, Initech 123 displays \texttt{Command refused} and moves nothing.
\end{Note}

\section{Moving a Formula: An Example}
A1 holds 10, B1 the formula \texttt{+A1*2}, and D3 the formula
\texttt{+B1+1}.  Move B1 to C5: select /Move with the cell pointer on B1,
press \kname{ENTER}, type \type{C5} and press \kname{ENTER}.  The moved
formula still uses A1, and the formula in D3 now refers to C5:

\LotExample{examples/move}

\section{Moving a Named Range: An Example}
B1..B2 is named DATA, and C1 adds it with \texttt{@SUM(DATA)}.  Move
B1..B2 to D1: the name moves with its cells, and the formula in C1, which
the control panel shows with the cells the name stands for, now adds
D1..D2.  A table of the names (\slashcmd{/Range Name Table} at F1) shows
the moved range:

\LotExample{examples/movename}
